\begin{examplebox}

	\textbf{\textcolor{CExample}{例1（基础）}}

	\textbf{\textcolor{CExample}{题目：}}
	在正方体$ABCD-A_1B_1C_1D_1$中，$O$为底面$ABCD$对角线$AC$与$BD$的交点。求证：$BD\perp A_1O$。

	\textbf{\textcolor{CExample}{解题步骤：}}
	\begin{description}[style=nextline, leftmargin=2.8em, labelsep=0.5em, itemsep=0.45em]
		\item[\textbf{第一步（确定射影）：}]
			\[
				A_1A\perp ABCD.
			\]
			\par\textbf{理由：} $A_1A\perp$底面$ABCD$，且$O$在底面内，故$A_1O$在底面内的射影为$AO$。

		\item[\textbf{第二步（证面内垂直）：}]
			\[
				BD\perp AO.
			\]
			\par\textbf{理由：}在正方形$ABCD$中，对角线互相垂直，故$BD\perp AO$。

		\item[\textbf{第三步（套用结论）：}]
			\[
				BD\perp A_1O.
			\]
			\par\textbf{理由：}由三垂线定理（正向），因为$BD\subset$底面$ABCD$且$BD\perp AO$（射影），所以$BD\perp A_1O$。
	\end{description}

	\textbf{\textcolor{CExample}{关键结论：}}
	\textbf{$\boxed{BD\perp A_1O}$}

	\medskip
	\textbf{\textcolor{CExample}{例2（稍变形）}}

	\textbf{\textcolor{CExample}{题目：}}
	在三棱锥$P-ABC$中，$PO\perp$底面$ABC$于$O$，$PA$在底面$ABC$内的射影为$AO$。若$BC\perp PA$，求证：$BC\perp AO$。

	\textbf{\textcolor{CExample}{解题步骤：}}
	\begin{description}[style=nextline, leftmargin=2.8em, labelsep=0.5em, itemsep=0.45em]
		\item[\textbf{第一步（线面垂直基础）：}]
			\[
				PO\perp ABC\;\Rightarrow\;PO\perp BC.
			\]
			\par\textbf{理由：}由$PO\perp$底面$ABC$及线面垂直定义得$PO\perp BC$。

		\item[\textbf{第二步（逆用定理）：}]
			\[
				BC\perp PA,\;PO\perp BC,\;PO\cap PA=P\;\Rightarrow\;BC\perp AO.
			\]
			\par\textbf{理由：}$BC$垂直于平面$PAO$内的两条相交直线$PO$和$PA$，故$BC\perp$平面$PAO$，从而$BC\perp AO$。也可直接由三垂线逆定理得到。

		\item[\textbf{第三步（结论）：}]
			\[
				BC\perp AO.
			\]
			\par\textbf{理由：}已证得$BC\perp AO$，结论成立。
	\end{description}

	\textbf{\textcolor{CExample}{关键结论：}}
	\textbf{$\boxed{BC\perp AO}$}

	\medskip
	\textbf{\textcolor{CExample}{例3（综合应用）}}

	\textbf{\textcolor{CExample}{题目：}}
	在正三棱柱$ABC-A_1B_1C_1$中，$D$为$AB$中点，$AA_1\perp$底面$ABC$。求证：$CD\perp A_1B$。

	\textbf{\textcolor{CExample}{解题步骤：}}
	\begin{description}[style=nextline, leftmargin=2.8em, labelsep=0.5em, itemsep=0.45em]
		\item[\textbf{第一步（确定射影）：}]
			\[
				AA_1\perp ABC.
			\]
			\par\textbf{理由：}$AA_1\perp$底面$ABC$，故$A_1B$在底面内的射影为$AB$。

		\item[\textbf{第二步（证面内垂直）：}]
			\[
				CD\perp AB.
			\]
			\par\textbf{理由：}在正三角形$ABC$中，$D$为$AB$中点，所以$CD\perp AB$（三线合一）。

		\item[\textbf{第三步（套用结论）：}]
			\[
				CD\perp A_1B.
			\]
			\par\textbf{理由：}由三垂线定理，$CD\perp AB$（射影）且$CD\subset$底面$ABC$，故$CD\perp A_1B$。
	\end{description}

	\textbf{\textcolor{CExample}{关键结论：}}
	\textbf{$\boxed{CD\perp A_1B}$}
\end{examplebox}
